\section{Polynomial differential equations from finite arithmetic}
\label{sec:ode}

The three-impossibilities packet asks whether the analytic infrastructure
for the Hermite identity and disk-area calculation can be organized around
a generic theorem for polynomial differential equations with explicit
majorants \cite{impossibilitiespacket}.  We give such a construction, together
with a finite algebraic method for proving its conserved quantities.  This
extends the earlier libraries of names, finite integration, and implicit
branches \cite{oldframework}.  The new construction produces a primitive
recursive approximation scheme and its finite error estimates; its
existence proof never invokes completeness of the real numbers.

\subsection{Finite polynomial data and a constructed local horizon}

A rational polynomial vector field in dimension $d\ge1$ is a finite table
\[
 F_i(x)=\sum_{\alpha}c_{i,\alpha}x^\alpha,
 \qquad x^\alpha=\prod_{j<d}x_j^{\alpha_j},\qquad
 |\alpha|=\sum_{j<d}\alpha_j.
\]
Rational arithmetic, exponent vectors, and all the tables are numerically
coded.  Polynomial evaluation is uniform primitive recursion on this
particular finite syntax; it is not evaluation of arbitrary primitive
recursive program codes.  Dimension, degree, coefficient sizes, and
precision may all vary.  We use the norm
$\lVert x\rVert=\max_{i<d}|x_i|$.

For rational $B\ge1$, compute
\begin{align}
 M&=\max\left(1,\max_{i<d}\sum_\alpha
                  |c_{i,\alpha}|B^{|\alpha|}\right),\label{eq:ode-M}\\
 L&=\max_{i<d}\sum_{|\alpha|\ge1}
                  |c_{i,\alpha}|\,|\alpha|B^{|\alpha|-1}.
                  \label{eq:ode-L}
\end{align}

\begin{lemma}[Finite polynomial bounds]
\label{lem:ode-polybound}
If $\lVert x\rVert,\lVert y\rVert\le B$, then
\[
 \lVert F(x)\rVert\le M,\qquad
 \lVert F(x)-F(y)\rVert\le L\lVert x-y\rVert.
\]
\end{lemma}
\begin{proof}
The first assertion is the triangle inequality applied to each monomial.
For the second, write a monomial as a product of $|\alpha|$ factors, with
coordinates repeated as prescribed by $\alpha$.  Replace the factors at
$x$ by the factors at $y$ one at a time.  The resulting telescoping identity
has $|\alpha|$ terms, each of absolute value at most
$B^{|\alpha|-1}\lVert x-y\rVert$.  Multiplication by coefficient absolute
values and finite summation give the result.  These are finite identities
and inequalities, including when dimension and degree vary.
\end{proof}

Given any rational initial vector $x_0$, the choices
\begin{equation}
 B=1+\lVert x_0\rVert,\qquad M,L\text{ from
 \eqref{eq:ode-M}--\eqref{eq:ode-L}},\qquad T=\frac1{2M}
 \label{eq:ode-localhorizon}
\end{equation}
construct a positive rational horizon.  For $N=2^n$ and $h=T/N$, define
\begin{equation}
 x_{n,0}=x_0,\qquad x_{n,j+1}=x_{n,j}+hF(x_{n,j})\quad(j<N).
 \label{eq:ode-euler}
\end{equation}
Induction gives
\begin{equation}
 \lVert x_{n,j}-x_0\rVert\le jhM\le TM=\tfrac12\quad(j\le N).
 \label{eq:ode-invariant}
\end{equation}
At each step the preceding estimate puts the current point in the box,
where Lemma~\ref{lem:ode-polybound} applies.  Thus both the local horizon
and its validity proof are produced from the finite input.  They are not
extra existence hypotheses.

The subsequent theorems also work on any larger rational horizon $T>0$
for which a uniform arithmetic proof establishes that the Euler nodes
remain in $\lVert x\rVert\le B$.  The finite inequality
$\lVert x_0\rVert+TM\le B$ is one sufficient certificate.  Concrete linear
and triangular systems below admit sharper invariant proofs.  Restricting
to meshes $2^{n+b}$, for a supplied natural $b$, is harmless: every bound
below remains true with $2^{-n}$ as an overestimate of the mesh factor.

\subsection{Discrete Gronwall with an entirely rational growth bound}

\begin{lemma}[Rational growth bound]
\label{lem:ode-growth}
For $a\in\Q$, $a\ge0$, and integer $k\ge1$,
\[
 \left(1+\frac ak\right)^k\le3^{\lceil a\rceil}.
\]
\end{lemma}
\begin{proof}
Let $A=\lceil a\rceil$.  Finite Bernoulli gives
$1+a/k\le1+A/k\le(1+1/k)^A$.  The binomial expansion gives
\[
 (1+1/k)^k\le\sum_{j=0}^k\frac1{j!}\le3,
\]
because $\binom{k}{j}\le k^j/j!$ and $j!\ge2^{j-1}$ for $j\ge2$;
the last step is a finite geometric sum.  Raise the Bernoulli inequality
to the $k$th power.  If $A=0$, then $a=0$ and the claim is immediate.
\end{proof}

\begin{lemma}[Finite Gronwall estimate]
\label{lem:ode-gronwall}
Suppose $e_0,\ldots,e_k$ and $\delta_0,\ldots,\delta_{k-1}$ are
nonnegative rationals and
$e_{j+1}\le r e_j+\delta_j$, where $r\ge1$.  Then
\[
 e_j\le r^j\left(e_0+\sum_{i<j}\delta_i\right)\quad(j\le k).
\]
In particular, if $r=1+a/k$, $k\ge1$, and $\delta_i\le\delta$, then
\begin{equation}
 e_j\le3^{\lceil a\rceil}(e_0+k\delta)\quad(j\le k).
 \label{eq:ode-gronwall}
\end{equation}
\end{lemma}
\begin{proof}
The first inequality is a finite induction.  In its step, multiply the
induction bound by $r$ and use
$\delta_j\le r^{j+1}\delta_j$.  For the second assertion, use
$r^j\le r^k$, Lemma~\ref{lem:ode-growth}, and the finite sum bound.
After rational coding and Booleanization of the bounded input conditions,
the invariant is quantifier-free.  It involves no analytic assertion.
\end{proof}

\subsection{A primitive recursive family of solution names}

Interpolate the nodes of \eqref{eq:ode-euler} by
\begin{equation}
 U_n(t)=x_{n,j}+(t-jh)F(x_{n,j})
 \quad(jh\le t\le(j+1)h),\qquad U_n(T)=x_{n,N}.
 \label{eq:ode-polygon}
\end{equation}
The containing cell is found by integer division.  At grid boundaries
the two formulas agree.  The invariant box is convex, so it contains the
whole polygon.  Every polygon is $M$-Lipschitz: split any rational time
interval at its finitely many mesh points and sum the bounds on the
constant slopes.

Compute
\begin{equation}
 A=\lceil LT\rceil,\qquad G=3^A,
 \qquad K=\max(1,\lceil GLMT^2\rceil).
 \label{eq:ode-constants}
\end{equation}

\begin{theorem}[Uniform finite refinement estimate]
\label{thm:ode-refine}
For every rational $0\le t\le T$,
\begin{equation}
 \lVert U_{n+1}(t)-U_n(t)\rVert
       \le\frac{GLMT^2}{2\,2^n}.
 \label{eq:ode-refine}
\end{equation}
Consequently, for all $m\ge n$,
\begin{equation}
 \lVert U_m(t)-U_n(t)\rVert\le K2^{-n}.
 \label{eq:ode-cauchy}
\end{equation}
\end{theorem}
\begin{proof}
Put $N=2^n$, $h=T/N$, and compare the two polygons at times $kh/2$,
$k\le2N$.  Write $y_k=U_n(kh/2)$, $z_k=U_{n+1}(kh/2)$ and
$e_k=\lVert y_k-z_k\rVert$.  At an even $k$ the coarse slope is
$F(y_k)$.  At an odd $k$, the point $y_k$ is a half-step from a coarse
node $x$, so $\lVert y_k-x\rVert\le Mh/2$.  Thus the coarse slope
differs from $F(y_k)$ by at most $LMh/2$.  In both cases,
\[
 e_{k+1}\le(1+Lh/2)e_k+LMh^2/4,
 \qquad e_0=0.
\]
Apply \eqref{eq:ode-gronwall} with $2N$ steps and $a=LT$.  It gives
$e_k\le GLMT^2/(2N)$ at every fine node.  On a fine interval both
polygons are affine, and their difference is the convex interpolation
of its endpoint values.  The same norm bound therefore holds throughout
the interval.  This proves \eqref{eq:ode-refine}.

For \eqref{eq:ode-cauchy}, telescope the finite list
$U_n(t),\ldots,U_m(t)$ and use
$\sum_{j=n}^{m-1}2^{-j-1}\le2^{-n}$.  No infinite sum or completeness
principle occurs.
\end{proof}

For integer $K\ge1$, let $\ell(K)$ be the least $r\le K$ such that
$K\le2^r$, found by bounded search.  With
\begin{equation}
 c=\ell(K)+3,\qquad X(p,t)=U_{p+c}(t),
 \label{eq:ode-name}
\end{equation}
the estimate \eqref{eq:ode-cauchy} gives
\[
 \lVert X(q,t)-X(p,t)\rVert\le2^{-p-3}\qquad(q\ge p).
\]
This is a uniform regular name, with spare precision.  The mathematical
object produced by the theorem is this primitive recursive rational
evaluator and these contracts, not an independently postulated real limit.

For completeness, the time input can itself be represented.  Given a fixed
regular time-name evaluator $\tau(p)$ representing a point in $[0,T]$,
clip its rational approximations to $[0,T]$.  Clipping does
not increase distance.  Put
$d_0=\ell(\max(1,\lceil M\rceil))+3$ and use
\[
 U_{p+c}\bigl(\operatorname{clip}(\tau(p+d_0))\bigr).
\]
Split the difference of two outputs into a same-time mesh error and a
same-mesh time error.  The former is at most $2^{-p-3}$, and the latter,
by $M$-Lipschitz continuity, is at most $2^{-p-3}+2^{-q-3}$ when $q\ge p$.
This gives the regular-name contract.  The evaluator for $\tau$ is fixed,
or supplied through a finite table with explicit demands.

\subsection{The solution contract and its derivative modulus}

The following definition states precisely which differential-equation
claim is proved.  In particular, it keeps the uniform error data visible.

\begin{definition}[Controlled solution name]
\label{def:controlled-flow}
For fixed finite data $F,x_0,B,T$ and rational $C\ge0$, a controlled
solution name is a primitive recursive rational evaluator $Y(p,t)$,
for rational $0\le t\le T$, such that:
\begin{enumerate}
\item it is regular in $p$, uniformly in rational $t$: for all $p,q$,
      $\lVert Y(p,t)-Y(q,t)\rVert\le2^{-p}+2^{-q}$;
\item $\lVert Y(p,t)\rVert\le B$ and
      $\lVert Y(p,0)-x_0\rVert\le2^{-p}$;
\item for rational $q,r\in[0,T]$,
\begin{equation}
 \begin{split}
 &\lVert Y(p,r)-Y(p,q)-(r-q)F(Y(p,q))\rVert\\
 &\hspace{30mm}\le C|r-q|^2+2^{-p}.
 \end{split}
 \label{eq:ode-controlled}
\end{equation}
\end{enumerate}
For a fixed evaluator these are first-order arithmetic formulas.
The construction of its code and the proof of its contracts are separate
from the assertion that it is a solution name.
\end{definition}

\begin{theorem}[Constructed controlled flow]
\label{thm:ode-existence}
The evaluator \eqref{eq:ode-name} is a controlled solution name with
$C=LM/2$.  Hence every rational polynomial vector field and rational
initial vector have a uniformly constructed positive local horizon and
a primitive recursive controlled solution there.
\end{theorem}
\begin{proof}
Regularity, the box bound and the initial value have already been proved.
Let $v=|r-q|$.  On the cell containing $s$, its left endpoint is within
$Mh$ of $U_n(s)$.  The latter is within $M|s-q|$ of $U_n(q)$.
Lemma~\ref{lem:ode-polybound} bounds the difference between the polygon
slope and $F(U_n(q))$ by $LM(h+|s-q|)$.

Split the interval between $q$ and $r$ at its finitely many mesh points.
On each resulting subinterval the slope difference is constant.  Evaluate
the affine upper bound at that subinterval's rational midpoint and multiply
by its length.  Summing these inequalities and the constant polygon slopes
yields
\begin{equation}
 \begin{split}
 &\lVert U_n(r)-U_n(q)-(r-q)F(U_n(q))\rVert\\
 &\hspace{30mm}\le LM\bigl(v^2/2+hv\bigr).
 \end{split}
 \label{eq:ode-localdefect}
\end{equation}
For $q\le r$, the midpoint-times-length sum of $s-q$ is exactly
$(r-q)^2/2$, by the difference-of-squares identity on each subinterval.
Thus this uses only finite rational arithmetic, with no general positivity
theorem for integrals.  Reversing the orientation proves
the same estimate for $r<q$.

Substitute $n=p+c$, use $v\le T$, and observe
$LMT^2\le K$ and $K2^{-c}\le1/8$.  The extra mesh term is at most
$2^{-p-3}$, so \eqref{eq:ode-controlled} follows with $C=LM/2$.
Finally \eqref{eq:ode-localhorizon} supplies all hypotheses for arbitrary
finite rational polynomial input.
\end{proof}

The name-level interpretation of \eqref{eq:ode-controlled} is
\[
 \lVert Y(r)-Y(q)-(r-q)F(Y(q))\rVert\le C|r-q|^2.
\]
For nonzero rational $r-q$, the difference-quotient error is at most
$C|r-q|$.  Thus
$|r-q|\le2^{-k}/\max(1,\lceil C\rceil)$ is a derivative modulus,
with one-sided derivatives at endpoints.  This interpretation can be
proved entirely as all-precision rational tests using
\eqref{eq:ode-controlled}; no exact real sign or equality decision is used.
The same contract gives a time-continuity modulus, since for $|r-q|\le1$
the increment is bounded by $(M+C)|r-q|+2^{-p}$ at precision $p$.

\subsection{Uniqueness proved at one finite precision}

\begin{theorem}[Quantitative uniqueness]
\label{thm:ode-unique}
Two controlled solution names for the same $F,x_0,B,T$, with respective
defect constants $C_1,C_2$, represent the same function on $[0,T]$.
There is an explicit primitive recursive schedule for the mesh and name
precision used to verify equality at any requested accuracy.
\end{theorem}
\begin{proof}
Fix a rational $t\in[0,T]$, a positive integer $N$, and a name precision
$p$.  Compare the rational approximations to the two candidates at
$jt/N$, $j\le N$.  If $e_j$ is their distance, the two instances of
\eqref{eq:ode-controlled} give
\[
 e_{j+1}\le(1+Lt/N)e_j+(C_1+C_2)(t/N)^2+2\,2^{-p},
 \qquad e_0\le2\,2^{-p}.
\]
The finite Gronwall estimate therefore gives
\begin{equation}
 e_N\le G\left((C_1+C_2)T^2/N+2(N+1)2^{-p}\right).
 \label{eq:ode-uniquefinite}
\end{equation}
For a requested accuracy $2^{-k}$ choose a dyadic $N$ at least
$4G(C_1+C_2)T^2 2^k$ (and at least one), then choose $p$ with
$2^p\ge16G(N+1)(1+LT)2^k$.  Both choices use bounded logarithms of explicit
integers.  Each term in \eqref{eq:ode-uniquefinite} is at most
$2^{-k-2}$.  The resulting arbitrarily accurate agreement, together with
regularity, is exactly equality of the two name families: a primitive
recursive common reindexing makes their rational outputs differ by at
most $2^{-k}$ for every $k$.  Regularity transfers this agreement to every
finer requested observation.

In the object-theory proof, $N,p,t,k$ are numerical parameters and the
induction is on the rational error table $e_j$.  It is not induction on
the assertion that two real functions are equal.
\end{proof}

The theorem deliberately specifies controlled candidates.  It does not
extract a primitive recursive modulus from qualitative differentiability
of an arbitrary external function.  Its useful feature is that the
constructed solution supplies the requisite modulus automatically.

A related finite comparison is useful when composing constructions.  If
two polynomial fields have a common Lipschitz bound $L$ and a coefficient
estimate proves $\lVert F(x)-\widetilde F(x)\rVert\le\varepsilon$ on the
box, their Euler trajectories with initial separation $\eta$ satisfy
\[
 e_{j+1}\le(1+Lh)e_j+h\varepsilon.
\]
Their represented flows therefore obey
\begin{equation}
 \lVert X_F(t)-X_{\widetilde F}(t)\rVert
       \le G(\eta+T\varepsilon).
 \label{eq:ode-perturbation}
\end{equation}
An additional step error $\delta_j$ adds
$G\sum_j\delta_j$ to the bound.  This gives a predetermined precision
budget for rounded steps or Cauchy-named coefficients, provided the
common domain and invariance contracts are supplied.  The unconditional
local theorem above already covers arbitrary finite rational polynomial
codes.

\subsection{Finite algebraic identities produce conserved quantities}

Let $H(x)=\sum_\alpha b_\alpha x^\alpha$ be a rational polynomial.
Its formal partial derivatives are finite coefficient tables.

\begin{theorem}[Polynomial conservation compiler]
\label{thm:ode-conservation}
If the coefficientwise polynomial identity
\begin{equation}
 \sum_{i<d}(\partial_iH)(x)F_i(x)=0
 \label{eq:ode-conservationcert}
\end{equation}
holds, then the constructed flow satisfies
$H(X(t))=H(x_0)$ for every represented time $t\in[0,T]$.
\end{theorem}
\begin{proof}
Put $D=\max(1,B,MT)$ and compute the rational
\[
 R_H=\sum_{|\alpha|\ge2}|b_\alpha|\,2^{|\alpha|}
                                  D^{|\alpha|-2}.
\]
The multivariate binomial expansion gives the finite remainder estimate
\begin{equation}
 |H(x+v)-H(x)-\nabla H(x)\cdot v|\le R_H\lVert v\rVert^2
 \label{eq:ode-polyremainder}
\end{equation}
when $\lVert x\rVert\le B$ and $\lVert v\rVert\le MT$.
Indeed, degree-zero and degree-one terms in $v$ are precisely the
subtracted terms.  In every remaining term retain two factors of $v$
and bound the other factors by $D$.  For a monomial of total degree $e$,
the sum of the relevant binomial coefficients is at most $2^e$.
Finite summation proves \eqref{eq:ode-polyremainder}.

Take $v=hF(x)$.  Identity \eqref{eq:ode-conservationcert} makes the
linear term zero, so a full Euler step changes $H$ by at most
$R_HM^2h^2$.  At time $t=jh+\theta$, $0\le\theta\le h$, telescope
the $j$ full steps and the final partial step.  Since
$jh^2+\theta^2\le h(jh+\theta)\le hT$,
\begin{equation}
 |H(U_n(t))-H(x_0)|\le R_HM^2T^2\,2^{-n}.
 \label{eq:ode-conservationerror}
\end{equation}
The polynomial $H$ has an explicit Lipschitz bound on the box, by the
proof of Lemma~\ref{lem:ode-polybound}.  Combining it with
\eqref{eq:ode-cauchy} and \eqref{eq:ode-conservationerror}, at an
arbitrary specified precision, proves equality of the represented values.
The represented-time extension uses the same continuity bound.
\end{proof}

The certificate \eqref{eq:ode-conservationcert} is checked by finite
polynomial multiplication and rational coefficient comparison.  The
passage from that certificate to an exact invariant uses only the
displayed finite error bounds.  This is a useful general separation:
algebra finds the identity, while one common arithmetic theorem transports
it to the represented flow.

If instead $\nabla H\cdot F=J$, retaining the linear term gives the
corresponding integral identity.  One can also add an accumulator coordinate
$z'=J(x)$ and apply the zero-derivative theorem to $H(x)-z$.
The accumulator's increments are finite Riemann sums with an explicitly
vanishing defect.  Thus its identification with the integral uses the
existing finite-sum integration interface, without invoking an existence
theorem for integrals of arbitrary functions.

\subsection{A common engine for the Hermite identity and circle calculus}

\paragraph{Exponential and oscillator names.}
For $x'=x$, $x(0)=1$, on any rational horizon $T>0$, finite Euler states
are bounded by
$(1+T/N)^j\le3^{\lceil T\rceil}$.  Thus a uniform invariant box is
available directly from Lemma~\ref{lem:ode-growth}, without imposing
the smaller local horizon.  At $T=1$ the endpoint approximants are
$(1+2^{-n})^{2^n}$.  The generic construction supplies a name, its
derivative contract, and its uniqueness property.  Existing exponential
names with the same controlled residual and initial value can be identified
by Theorem~\ref{thm:ode-unique}.

For the oscillator
\[
 c'=-s,\qquad s'=c,\qquad(c(0),s(0))=(1,0),
\]
the sup norm of a step grows by at most $1+h$, so the same invariant
box works.  The coefficient identity
\[
 \partial_c(c^2+s^2)(-s)+\partial_s(c^2+s^2)c
       =-2cs+2sc=0
\]
gives $c(t)^2+s(t)^2=1$ through Theorem~\ref{thm:ode-conservation}.
Once the standard sine and cosine names satisfy the controlled derivative
contract, one uniqueness theorem identifies both constructions.

\paragraph{The finite Hermite antiderivative identity.}
Let $f\in\Q[t]$ have degree at most $m$, and compute
\[
 Q=\sum_{j=0}^{m}f^{(j)}.
\]
Formal differentiation is a finite coefficient operation, and telescoping
gives $Q-Q'=f$.  Consider the polynomial system
\begin{equation}
 \tau'=1,\qquad E'=-E,\qquad J'=E f(\tau),
 \qquad (\tau(0),E(0),J(0))=(0,1,0).
 \label{eq:hermite-flow}
\end{equation}
For any supplied rational horizon $T>0$, begin with meshes satisfying
$h\le1$, accomplished by a bounded logarithm.  The Euler invariant is
explicit: $0\le\tau_j\le T$, $0\le E_j\le1$, and
\[
 |J_j|\le jh F_0\le TF_0,\qquad
 F_0=\sum_{k\le m}|f_k|\max(1,T)^k.
\]
The bound on $E_j$ follows from $E_{j+1}=(1-h)E_j$; the bound on $J_j$
is a finite induction.  Therefore the generic flow theorem applies on
the entire supplied horizon.

The polynomial $H(\tau,E,J)=E Q(\tau)+J$ satisfies
\[
 \nabla H\cdot(1,-E,Ef(\tau))
   =E\bigl(Q'(\tau)-Q(\tau)+f(\tau)\bigr)=0.
\]
The conservation theorem proves
\begin{equation}
 J(t)=Q(0)-E(t)Q(t).
 \label{eq:hermite-flowidentity}
\end{equation}
Here $\tau(t)=t$, already exactly true of the Euler polygons, and $E$
is the controlled solution of $E'=-E$, $E(0)=1$.  The accumulator $J$
is the integral name of $E(t)f(t)$: sum its controlled increments over
a rational mesh.  At finite table precision $p$, the total defect is
bounded by $C_JT^2/N+(N+1)2^{-p}$, where $C_J$ is its controlled-increment
constant.  First choose $N$ to make the mesh term small, then choose $p$
to make the second term small.  The product-evaluation error has an
explicit coefficient bound and is included by increasing $p$ further.
Thus the sums converge with an explicit modulus to $J(t)$.
Identifying $E$ with the existing $e^{-t}$ name by uniqueness gives the
usual Hermite identity
\begin{equation}
 \int_0^t f(u)e^{-u}\,du=Q(0)-e^{-t}Q(t).
 \label{eq:hermite-identity}
\end{equation}
The integral notation in \eqref{eq:hermite-identity} denotes this
explicit finite-sum name.  Formula \eqref{eq:hermite-flowidentity} was
proved first, using only finite polynomial identities and Euler estimates.
The divisibility arguments and the final transcendence contradiction in
the original application remain separate arithmetic steps.

\paragraph{Circle-area accumulator.}
Add coordinates $\tau,A$ to the oscillator with
$\tau'=1$ and $A'=(c^2+s^2)/2$, initially zero.  The same finite
growth bound supplies an Euler box for $c,s$, and
$|A_j|\le T\,3^{2\lceil T\rceil}$ is a sufficient accumulator bound.
The two polynomial certificates are
\[
 \nabla(c^2+s^2)\cdot F=0,
 \qquad
 \nabla\bigl(A-\tau(c^2+s^2)/2\bigr)\cdot F=0.
\]
Consequently $c^2+s^2=1$ and $A(t)=t/2$.  Since
$(c s'-s c')/2=(c^2+s^2)/2$, $A$ is the signed sector-integral
accumulator.  The identification of a full period with $2\pi$ and the
geometric interpretation of that integral as disk area are still supplied
by the relevant geometric argument.  What has been resolved is the
common analytic construction and its exact polynomial identities, not
every step of those applications.

\subsection{The iPRA induction audit}

All algorithms here are primitive recursive uniformly in their finite
inputs.  The large value $2^{p+c}$ is harmless for primitive recursiveness;
the theorem makes no feasibility claim.  The required inductions have
explicit finite invariants: the Euler box estimate, the Gronwall bound
with partial error sum, the paired fine/coarse-node comparison, the
geometric telescoping bound, and the accumulated conservation defect.
Rational inequalities are reduced to integer inequalities by positive
denominators.  Finite coordinate, monomial, and mesh quantifiers are
Booleanized before induction.

For equality or differentiability of names, first fix a numerical
precision challenge.  The displayed bounds compute a sufficient mesh
and finite evaluation table.  Prove the resulting finite arithmetic
assertion, and only then introduce the challenge universally.  This is
the point at which one avoids induction on an analytic formula.

The mathematical construction therefore has a concrete path to
$\mathsf{Prf}\,\mathsf{PRAAx}$ proofs: code the finite algorithms, prove
their defining equations, and implement those bounded invariants in the
existing object calculus.  The complete deep Lean derivations are not
claimed here.  Ambient Lean correctness of a rational Euler algorithm
and an object-language iPRA proof remain distinct deliverables.
